% Recuperación documental para el artículo V; RESULTADO_RECUPERADO.
% Propietario: /Users/ruben/Documents/New project/output/REVISION_EDITORIAL_HMT_MD_20260905/01_FUENTES/INTEGRAL/tree/New project/output/ENSAMBLADO_CAUSAL_RESTAURADO_HMT_MD_20260905/fuente/manuscrito/integracion_83/parte_iv/body/B011_ch53_trayectorias_cuanticas_6b08e0e532c9_04ab_postulados_cuanticos_genealogicos_rev2.tex
% Líneas de origen: 4--95, 178--206.
% Sólo se adaptan las etiquetas de referencia y, cuando procede, el nivel de títulos.
% La matriz de recuperación conserva la correspondencia y las dependencias.
\section{Reconstruction des principes quantiques à partir de l’état généalogique}
\label{v:rec:sec:postulados-cuanticos-genealogicos-rev2}
\index{mécanique quantique!reconstruction généalogique}
\index{état!réalisation de Hilbert}

La formulation de Hilbert organise la mécanique quantique au moyen d’états,
d’observables autoadjoints, d’une évolution unitaire, d’une composition tensorielle et d’une
règle probabiliste pour les résultats~\cite{vonNeumann1932}. HMT--MD
construit une couche antérieure à cette représentation : une section compatible,
son chemin observable, la mémoire du transport et le lecteur qui exprime un
résultat. L’espace de Hilbert n’est pas remplacé ; il apparaît comme codomaine d’une
réalisation linéaire déclarée.

La suite typée est
\[
 \boxed{\begin{aligned}
 \text{section généalogique}
 &\longrightarrow \text{réalisation linéaire}
 \longrightarrow \text{phase et superposition}\\
 &\longrightarrow \text{observable et évolution}
 \longrightarrow \text{probabilité et mesure}
 \end{aligned}}
\]
Chaque flèche conserve son domaine. En particulier, le vecteur d’état ne
remplace pas la section TPK, de même que la valeur réelle ne remplace pas
la généalogie numérique qui la produit.

\subsection{État généalogique et réalisation de Hilbert}
\label{v:rec:realizacion_evolucion:section:02}

\begin{definition}[État généalogique réalisé]
\label{v:rec:realizacion_evolucion:statement:01}
Un état généalogique est un quintuplet
\begin{equation}
 \mathfrak G=(x_\infty,\Gamma,\mathcal F,\mathcal M,\mathcal R),
 \label{v:rec:eq:estado-genealogico-realizado-rev2}
\end{equation}
où \(x_\infty\in\Omega_\infty\) est une section compatible,
\(\Gamma\) son chemin observable, \(\mathcal F\) la collection des réalisations,
\(\mathcal M\) la mémoire des relèvements et \(\mathcal R\) le lecteur
déclaré. Soit \(\mathsf P_{\rm TPK}\) la catégorie dont les objets sont les
sections TPK préparées, finies ou coinductives, et dont les morphismes sont les
chemins enregistrés compatibles avec leurs données de bord. En particulier, la
restriction de profondeur est un morphisme
\[
 r_{n+1,n}:x_{n+1}\longrightarrow x_n.
\]
Soit \(\mathsf{Hilb}_{\mathbb R,J}\) la catégorie dont les objets
sont des paires \((\mathcal H_{\mathbb R},J)\), avec \(\mathcal H_{\mathbb R}\)
un espace de Hilbert réel et \(J\) orthogonal, \(J^2=-I\), et dont les morphismes sont
des applications linéaires réelles bornées \(T\) qui entrelacent les structures
complexes, \(TJ_1=J_2T\). Une réalisation quantique est un foncteur
\[
 \mathcal Q:\mathsf P_{\rm TPK}\longrightarrow
 \mathsf{Hilb}_{\mathbb R,J},
\]
et une \emph{réalisation pointée} est la paire \((\mathcal Q,\Psi)\) qui associe
à chaque section préparée \(x\) un vecteur non nul
\(\Psi_x\in\mathcal Q(x)\), compatible avec les restrictions du système
projectif.
\end{definition}

Avec la convention de linéarité dans le second argument, la forme hermitienne associée
s’écrit
\[
 \langle u,v\rangle_{\mathbb C}
 =\langle u,v\rangle_{\mathbb R}
  -i\langle u,J_E v\rangle_{\mathbb R}.
\]
Choisir le système de coordonnées \(J_E\leftrightarrow i\) réalise la structure complexe qui,
dans le noyau formel et son développement, est construite comme racine interne de \(-I\). Pour des préfixes de
profondeur \(n\), posons
\(\mathcal H_n:=\mathcal Q(x_n)\) et
\[
 p_{n+1,n}:=\mathcal Q(r_{n+1,n}):
 \mathcal H_{n+1}\longrightarrow\mathcal H_n.
\]
La compatibilité entre restriction et réalisation exige
\begin{equation}
 p_{n+1,n}\Psi_{n+1}(x_{n+1})=\Psi_n(x_n).
 \label{v:rec:eq:compatibilidad-realizacion-hilbert-rev2}
\end{equation}
Si \(r_{\infty,n}:x_\infty\to x_n\) est la restriction coinductive et
\(p_n:=\mathcal Q(r_{\infty,n}):\mathcal H_\infty\to\mathcal H_n\), une
limite \(\mathcal H_\infty\) réalise la section complète lorsque
\(p_n\Psi_\infty(x_\infty)=\Psi_n(x_n)\) à chaque profondeur.

\begin{principle}[État quantique réalisé]
\label{v:rec:realizacion_evolucion:statement:02}
L’état pur associé à une section est représenté par le rayon du vecteur
non nul \(\Psi=\Psi_\infty(x_\infty)\) ; une préparation statistique est
représentée par un opérateur densité positif et de trace un sur la même
réalisation. L’identification de phase globale agit dans la représentation :
la section et son registre restent dans le domaine généalogique.
\end{principle}

\subsection{Évolution, Hamiltonien et équation de Schrödinger}
\label{v:rec:realizacion_evolucion:section:03}

L’évolution généalogique préserve la composition des transports. Sa
réalisation satisfait
\[
 U(t_3,t_2)U(t_2,t_1)=U(t_3,t_1),
 \qquad U(t,t)=I.
\]
Lorsqu’elle préserve le produit scalaire et satisfait
\(U(t)J_E=J_EU(t)\), elle est orthogonale dans le système de coordonnées réel et unitaire dans le système de coordonnées
complexe. Pour un groupe unitaire à un paramètre fortement continu, le
théorème de Stone fournit un générateur autoadjoint et \(J_E\)-linéaire
\(H\)~\cite{Stone1932} :
\[
 U(t)=\exp\!\left(-J_E H t/\hbar_{\rm int}\right).
\]
Sur le domaine \(D(H)\), la dérivation temporelle donne
\begin{equation}
 \hbar_{\rm int}J_E\dot\Psi(t)=H\Psi(t),
 \label{v:rec:eq:schrodinger-genealogico-rev2}
\end{equation}
et le système de coordonnées \(J_E\leftrightarrow i\) produit
\(i\hbar_{\rm int}\dot\Psi=H\Psi\).

Le calendrier \(C_{108}\) réalise cette correspondance de manière finie :
le Hamiltonien autoadjoint construit dans sa source a pour
exponentielle, pendant \(t_0\), le décalage unitaire d’une itération. Le
groupe continu requiert en outre la continuité forte ; il ne se déduit pas
seulement de la périodicité du calendrier.
